\chapter{Requirements Analysis}
\label{ch:req}

This chapter turns the challenge statement into a requirements baseline. The challenge
names five qualities --- structured, inclusive, interactive, personalized, effective ---
and four roles the companion may adopt. Those are design intentions, not testable
requirements, so the work of this chapter is to say what each one means operationally,
and to fix the constraints under which it must be achieved.

\section{The four roles, read precisely}

Each role in the challenge statement implies a different relationship between the AI and
the group. Reading them separately was productive, because the naive design satisfies one
and violates the others.

\begin{table}[htbp]
\centering
\caption{The four roles named by the challenge, and what each rules out.}
\label{tab:roles}
\footnotesize
\begin{tabularx}{\textwidth}{@{}p{2.2cm}X X@{}}
\toprule
\textbf{Role} & \textbf{What HALQA does} & \textbf{What it forbids} \\
\midrule
Mediator & Opens a timed window, invites quiet members, routes the floor so
           participation stays balanced (S-B-U \emph{Share} ring)
         & Letting the fastest typist own the session; ranking learners publicly \\
Teammate & Prepares a Quick Take silently during the peer window, then releases it as one
            voice among several, with citations
         & Speaking first; being the only source of an answer \\
Tutor    & Hint ladder L1--L4, Socratic prompts, register-matched explanation, style
            adapted by bandit
         & Handing over a full solution before the learner has worked \\
Facilitator & Agenda, capsule, S-B-U health, \code{needs\_human} escalation to a
               mentor
         & Replacing mentor judgement; auto-answering on a mentor's behalf \\
\bottomrule
\end{tabularx}
\end{table}

The interesting case is \emph{teammate}. A teammate who always has the answer stops
being a teammate and becomes an oracle. We resolved this by separating \emph{preparation}
from \emph{release}: Nova is allowed to know the answer during the peer window, and is
forbidden to reveal it until the window closes. This single distinction is what makes the
Circle-First Ask protocol work, and it is enforced server-side rather than in the prompt
(\S\ref{sec:cfas}).

\section{The five qualities, made operational}

\begin{table}[htbp]
\centering
\caption{Qualities named by the challenge, translated into measurable system properties.}
\label{tab:qualities}
\footnotesize
\begin{tabularx}{\textwidth}{@{}p{1.9cm}p{4.5cm}X@{}}
\toprule
\textbf{Quality} & \textbf{Operational meaning} & \textbf{Mechanism} \\
\midrule
Structured & A session has an agenda, a sequence of question states, and a defined end
           & Session state machine \code{OPEN}$\to$\code{NOVA\_THINKING}$\to$\code{RESOLVED};
             agenda templates; revision capsule \\
Inclusive  & Participation is balanced, and shy learners are not penalized for being slow
           & Anonymous questions, private ``me too'', log-capped credits so volume cannot
             dominate, S-B-U \emph{Share} ring, private coach nudges \\
Interactive & Every learner acts between AI turns
            & 45--60\,s peer window, consensus polls, Whisper relay, quiz battles,
              solo ``teach Nova'' mode \\
Personalized & Explanation style, level, language and hint eagerness differ per learner
             & DNA passport (4 items), six-dimensional style weights updated by a bandit,
               four registers including Tunisian Darija \\
Effective  & The circle reaches a verified answer, and remembers it
           & Citations on every grounded answer, Auditor verification, capsule at session
             end, quiz bank derived from the session \\
\bottomrule
\end{tabularx}
\end{table}

\section{Functional requirements}

Requirements are grouped by the surface they expose. Priority follows MoSCoW; the
``Evidence'' column points to the section that demonstrates the requirement is met rather
than merely intended.

\begin{longtable}{@{}p{1.15cm}p{6.5cm}p{1.35cm}p{4.4cm}@{}}
\caption{Functional requirements.}\label{tab:fr}\\
\toprule
\textbf{ID} & \textbf{Requirement} & \textbf{Prio.} & \textbf{Evidence} \\
\midrule
\endfirsthead
\multicolumn{4}{@{}l}{\footnotesize\sffamily\color{ink2}Table \ref{tab:fr} continued}\\
\toprule
\textbf{ID} & \textbf{Requirement} & \textbf{Prio.} & \textbf{Evidence} \\
\midrule
\endhead
\bottomrule
\endfoot
FR-01 & A learner can create a circle and receive a join code & Must & \S\ref{sec:circles} \\
FR-02 & A learner can join by code or by QR, and the lobby lists members & Must & \S\ref{sec:circles} \\
FR-03 & A question opens a timed peer window of 30--90\,s (default 45) & Must & \S\ref{sec:cfas} \\
FR-04 & Peers answer inside the window; answers are attributed or anonymous & Must & \S\ref{sec:cfas} \\
FR-05 & The AI answer is prepared during the window and hidden until it closes & Must & \S\ref{sec:cfas} \\
FR-06 & Peer answers are verified before consensus is declared & Must & \S\ref{sec:auditor} \\
FR-07 & A four-item diagnostic produces a passport that changes later turns & Must & \S\ref{sec:dna} \\
FR-08 & A hint ladder L1--L4 is served, with L4 locked until L1--L3 are used & Must & \S\ref{sec:hints} \\
FR-09 & Every grounded answer cites knowledge-base chunk identifiers & Must & \S\ref{sec:rag} \\
FR-10 & A contested claim can be put to a poll, and votes are counted server-side & Should & \S\ref{sec:auditor} \\
FR-11 & A revision capsule of at most three bullets is produced at session end & Must & \S\ref{sec:capsule} \\
FR-12 & A session can end idempotently, so a remount does not double-credit & Must & \S\ref{sec:credits} \\
FR-13 & Solo mode runs the same ladder with the learner teaching the AI & Should & \S\ref{sec:solo} \\
FR-14 & The Whisper relay compares a paraphrase chain against a claim checklist & Should & \S\ref{sec:whisper} \\
FR-15 & A mentor can observe without earning credits or taking turns & Should & \S\ref{sec:mentor} \\
FR-16 & A facilitator board shows per-circle S-B-U rings and \code{needs\_human} & Should & \S\ref{sec:sbu} \\
FR-17 & Circles can be matched deterministically from passports & Could & \S\ref{sec:matcher} \\
FR-18 & A dashboard aggregates progress without exposing individual ranking & Should & \S\ref{sec:dashboard} \\
FR-19 & \code{/metrics} exposes tier counts, L-share, latencies and cache-hit rate & Must & \S\ref{sec:obs} \\
FR-20 & The interface is available in Arabic, Darija, French and English & Must & \S\ref{sec:i18n} \\
FR-21 & An offline pack serves content when the provider is unavailable & Should & \S\ref{sec:degradation} \\
FR-22 & A learner can export their data and request deletion & Must & \S\ref{sec:datasubject} \\
\end{longtable}

\section{Non-functional requirements}

\begin{table}[htbp]
\centering
\caption{Non-functional requirements and the budget they set.}
\label{tab:nfr}
\footnotesize
\begin{tabularx}{\textwidth}{@{}p{1.1cm}p{3.3cm}X p{3.3cm}@{}}
\toprule
\textbf{ID} & \textbf{Property} & \textbf{Target} & \textbf{Verification} \\
\midrule
NFR-01 & Time to first paint of a section skeleton & $<200$\,ms, zero model calls
       & Router config, unit test on budget \\
NFR-02 & Compact-tier completion & $<4$\,s & \code{config/router.yaml}, \code{/metrics} p50/p95 \\
NFR-03 & Larger-tier completion & $<12$\,s & \code{config/router.yaml}, \code{/metrics} \\
NFR-04 & Larger-model share of calls & $<0.25$ & \code{lShare} from the call ledger \\
NFR-05 & Per-session quota & $\le 20$ C, $\le 5$ L & Governor ledger, atomic reserve \\
NFR-06 & Per-circle daily ceiling & 400 calls & \code{CIRCLE\_DAILY\_CALLS} \\
NFR-07 & Peer state poll latency & p95 $<50$\,ms & \code{load:poller} SLO gate \\
NFR-08 & Accessibility & WCAG 2.1 AA, 44\,px targets, keyboard-complete
       & axe-core gate in CI \\
NFR-09 & Privacy & No name, photo, gender or region column & Schema gate test \\
NFR-10 & Security & Default-deny authz; strict CSP; no secrets in repo & Guard, CSP test, secret scan \\
NFR-11 & Availability under failure & Usable in all four degradation states
       & Chaos toggle plus smoke \\
NFR-12 & Reproducibility & Clean clone builds and tests from tracked files only & \code{clean-clone} CI job \\
\bottomrule
\end{tabularx}
\end{table}

\section{Constraints}

The constraints below are properties of the environment, and several of them shaped the
architecture more than any functional requirement did.

\paragraph{Constrained inference.} The provided inference platform runs without GPU
acceleration and access is shared. Concretely: latency is variable and occasionally
seconds-scale, throughput is not ours to scale, and the larger model is a scarce
resource that may be throttled or unavailable. A design that assumes a fast, elastic
provider would fail on the first shared session.

\paragraph{Shared-device, mixed-network context.} Target users are students in
universities, training centres and IEEE student branches. Sessions may run on shared
laptops and campus Wi-Fi that drops. The client is therefore a progressive web app with
an installable shell, an offline pack, and no assumption of a stable connection during a
session.

\paragraph{Four registers, one of them code-switched.} The challenge requires
multilingual behaviour including Arabic, French and English; the Tunisian context adds
Darija written in Arabic script and, crucially, Arabizi --- Latin-script transliteration
with digits standing in for letters (\code{3} for \emph{ain}, \code{7} for \emph{ha},
\code{9} for \emph{qaf}). Arabizi is not a dialect variant of a written language; it is a
different encoding that breaks naive tokenization, so it needs its own normalization
before retrieval or classification (\S\ref{sec:arabizi}).

\paragraph{Privacy law and ethics.} Learners may be minors. The system therefore collects
no direct identifiers, and the constraint is structural: the persona table has no column
in which a name, photo, gender or region could be stored.

\section{Use cases}

The principal actors are the \kw{learner}, the \kw{mentor} (observer, no turns), the
\kw{facilitator} (circle owner, reads the board) and the \kw{administrator}. Table
\ref{tab:usecases} lists the load-bearing use cases; the two most important are expanded
as scenarios afterwards, because they carry most of the design.

\begin{table}[htbp]
\centering
\caption{Principal use cases by actor.}
\label{tab:usecases}
\footnotesize
\begin{tabularx}{\textwidth}{@{}p{1.9cm}X p{3.0cm}@{}}
\toprule
\textbf{Actor} & \textbf{Use case} & \textbf{Requirement} \\
\midrule
Learner & Complete the DNA diagnostic and receive a passport & FR-07 \\
Learner & Create or join a circle, then enter the lobby & FR-01, FR-02 \\
Learner & Ask a question and wait through the peer window & FR-03--FR-05 \\
Learner & Answer a peer's question; earn quality-weighted credits & FR-04 \\
Learner & Request hints through the ladder & FR-08 \\
Learner & Vote on a contested claim & FR-10 \\
Learner & End the session and read the capsule & FR-11, FR-12 \\
Learner & Practise alone by teaching Nova & FR-13 \\
Learner & Play the Whisper relay and repair a drift & FR-14 \\
Mentor & Observe a circle and leave a note and a brief & FR-15 \\
Facilitator & Read the board, spot a circle needing a human & FR-16 \\
Facilitator & Match learners into balanced circles & FR-17 \\
Admin & Read \code{/metrics} and operate the chaos toggle & FR-19 \\
\bottomrule
\end{tabularx}
\end{table}

\subsection{Scenario A: the Circle-First Ask}

A learner posts \emph{``Why does binary search need a sorted list?''}. The server
creates the question in state \code{OPEN} with a 45-second window and returns
immediately. Peers see the question and answer; one writes that binary search works on
any array but is slower on unsorted ones, which is wrong. Nova, in parallel, retrieves
chunks, checks the semantic cache, and prepares a Quick Take grounded in five chunk
citations. When the window closes, the expiry worker settles the question: the Auditor
compares peer answers against the claim checklist and marks the incorrect one
\emph{doubtful}, which opens a poll rather than a counter-assertion. The learner then
receives, in order, the verified peer answer, the poll, and the citation-backed
explanation --- each stage announced, none of them a frozen screen.

\subsection{Scenario B: degradation under quota exhaustion}

A workshop runs twelve circles at once. The larger-model budget is exhausted
mid-session. The governor refuses further L reservations atomically; the router
downgrades the request to the compact tier; the degradation machine moves
\code{orange}$\to$\code{red} and emits a banner, a telemetry event and a metric gauge.
Learners keep receiving compact-tier answers with simplified labels, and the capsule at
session end is drafted at the compact tier without its polishing pass. Nothing in the
session fails; what changes is the ceiling on prose quality, which the interface states
honestly.

\section{Traceability}

Each functional requirement in Table~\ref{tab:fr} names the section that demonstrates it,
and each non-functional target in Table~\ref{tab:nfr} names the artefact that measures
it. The inverse mapping --- code to requirement --- is maintained by the test suite:
schema-gate, adapter-isolation, CSP and quota tests are written against requirements
rather than against implementation details, so a refactor that breaks a guarantee fails a
test whose name states the guarantee (\S\ref{ch:testing}).
